\chapter{Contenu de la réponse}

\noindent\fcolorbox{orangeECSR}{orangeECSR!7}{\parbox{0.94\textwidth}{\textbf{Statut de ce chapitre : récit rédigé à titre provisoire avant l'action.} Les verbes au passé sont une mise en forme anticipée, demandée pour préparer la version finale du mémoire ; ils ne constituent pas une attestation de réalisation. Après le 4 décembre 2026, les horaires, activités, réactions, observations et évaluations devront être confrontés aux faits et corrigés.}}

\section{Contenu et déroulement de l'action}
\subsection{Accueil, présentation et ancrage dans le métier de plombier}
Nous avons introduit notre séance par le titre de notre support : \og Smartphone au volant : le piège des chantiers \fg{}. Notre intention était de placer les stagiaires dans une situation professionnelle concrète et de montrer que le risque routier ne concerne pas uniquement le trajet domicile--formation : il peut aussi survenir lors des déplacements entre une entreprise, un fournisseur, un client et un chantier.

Après la présentation de notre binôme et le rappel du cadre de cette sensibilisation, nous avons évoqué quatre situations familières ou plausibles dans le métier : l'appel urgent d'un client, la recherche d'une adresse sur le GPS, le retard pris sur un chantier et l'envoi d'une photographie à un responsable. Nous avons posé une question d'accroche : \og Lorsqu'un client vous appelle pour une urgence alors que vous conduisez, devez-vous répondre immédiatement ? \fg{} Cette question devait susciter les premières réactions, sans jugement et sans chercher d'emblée une réponse unique.

Nous avons ensuite annoncé les quatre objectifs affichés dans le diaporama : identifier les risques liés au téléphone et aux distracteurs ; comprendre leurs conséquences sur la conduite ; connaître les règles et sanctions ; adopter des réflexes sûrs pour l'activité professionnelle. Nous avons souligné que la perte d'attention peut avoir des conséquences humaines et matérielles, mais aussi compromettre l'exercice d'un métier nécessitant des déplacements. \emph{La formulation du support selon laquelle la perte du permis entraînerait systématiquement la perte de l'employabilité a été traitée comme un message de vigilance, et non comme une règle applicable à toute situation individuelle.}

\subsection{Identifier les quatre dimensions de la distraction}
Nous avons organisé un apport dialogué autour des quatre dimensions de la distraction présentées dans les diapositives 8 à 12 :
\begin{description}[leftmargin=3.0cm,style=nextline]
\item[Visuelle] Détourner le regard de l'environnement routier, par exemple pour consulter un message, une notification ou l'écran du GPS.
\item[Manuelle] Mobiliser une main pour composer un numéro, manipuler un téléphone ou régler un équipement au lieu de garder les gestes nécessaires à la conduite.
\item[Auditive] Voir son attention sollicitée par une sonnerie, une conversation ou un signal sonore, notamment dans une situation d'urgence professionnelle.
\item[Cognitive] Consacrer ses ressources mentales à un échange stressant, à un problème de chantier ou à l'organisation de la prochaine intervention plutôt qu'à l'analyse de la circulation.
\end{description}
Nous avons insisté sur le fait qu'une même situation pouvait cumuler plusieurs dimensions : lire un SMS associe, notamment, distraction visuelle, manuelle et cognitive. Le cas du téléphone en mains libres a servi à distinguer la conformité de certains dispositifs avec la réglementation de l'absence de risque : une conversation peut encore mobiliser l'attention du conducteur.

\subsection{Deux activités vidéo sur les limites de l'attention}
Pour rendre les notions moins abstraites, nous avons utilisé deux exercices d'observation figurant dans notre support. Le premier consistait à compter des ballons selon une consigne donnée ; le second demandait de compter des passes dans une vidéo. Les participants devaient conserver leur réponse jusqu'à la fin du visionnage, puis discuter de ce qu'ils avaient vu et de ce qui pouvait leur avoir échappé.

Notre objectif n'était pas de mesurer médicalement l'attention des stagiaires, mais de faire comprendre qu'une tâche ciblée peut mobiliser fortement les ressources cognitives et réduire la perception d'autres événements. Nous avons relié ce phénomène au fait de rechercher une adresse, de répondre à un client ou de consulter un écran pendant la conduite. Les vidéos étaient des supports de discussion ; leurs résultats individuels ne devaient pas donner lieu à un classement des participants.

\subsection{Chiffres, distance parcourue et conséquences professionnelles}
Nous avons poursuivi par un éclairage sur la sécurité routière, à partir des diapositives consacrées aux données de l'ONISR et aux risques des distracteurs. Le support présentait notamment des chiffres de mortalité et de blessés en 2025, ainsi que des données de 2024 concernant le défaut d'attention et l'usage du téléphone chez les conducteurs de véhicules utilitaires. Nous avons présenté ces données comme des éléments de contextualisation et non comme des résultats issus du groupe de stagiaires. \emph{Les valeurs précises, leurs définitions et leurs sources devront être revérifiées dans la version définitive remise au jury.}

Nous avons ensuite proposé un calcul simple pour illustrer l'effet d'un regard détourné pendant cinq secondes. À vitesse constante, la distance parcourue est égale à la vitesse exprimée en mètres par seconde multipliée par le temps. Nous avons présenté les ordres de grandeur figurant dans le support : environ 70 mètres à 50 km/h, 110 mètres à 80 km/h, 125 mètres à 90 km/h et 180 mètres à 130 km/h. La discussion ne visait pas à prétendre qu'un conducteur ne perçoit nécessairement rien pendant cinq secondes ; elle montrait qu'un véhicule continue de se déplacer pendant que son conducteur réduit sa prise d'informations.

Les exemples d'un piéton traversant, d'un véhicule freinant ou d'une intervention sur chaussée ont permis de faire le lien avec les risques concrets. Nous avons rapproché ces situations de la pression professionnelle : désir de satisfaire le client, retard sur le planning, consignes du responsable ou urgence apparente.

\subsection{Cadre réglementaire et réflexes de prévention}
Nous avons consacré une séquence aux règles applicables à l'usage du téléphone au volant, en prenant pour référence l'article R.412-6-1 du Code de la route. Le support rappelait notamment l'interdiction de tenir un téléphone en main et celle de porter à l'oreille certains dispositifs émettant du son, sous réserve des exceptions réglementaires. Les sanctions mises en avant étaient une amende forfaitaire de 135 euros et le retrait de trois points pour l'infraction correspondante ; des conséquences supplémentaires peuvent exister dans certaines situations prévues par les textes. Nous avons distingué la sanction juridique du risque d'accident, qui demeure l'enjeu central de la sensibilisation.

Nous avons terminé cette partie par la recherche de mesures préventives applicables au métier : programmer son GPS avant de partir, activer si possible le mode conduite, placer le téléphone hors de portée, prévenir les clients de son indisponibilité temporaire, et s'arrêter sur un emplacement autorisé et sécurisé avant de rappeler ou de consulter un message. Nous avons insisté sur le rôle de l'organisation du travail : un appel considéré comme urgent ne doit pas imposer au conducteur de traiter une sollicitation en mouvement.

\subsection{Activité extérieure : « Même parcours, attention différente »}
Pour prolonger les apports théoriques, nous avons organisé une expérimentation exclusivement \textbf{pédestre}, sur un parcours balisé de huit à douze mètres. Ce parcours comportait quatre ou cinq plots à contourner, une zone d'arrêt, un objet à toucher ou récupérer, puis une ligne d'arrivée. La configuration permettait de reproduire une tâche simple, sans circulation de véhicule et sans recherche de vitesse.

L'effectif de référence étant de quinze stagiaires, nous avons retenu des passages par volontaires, les autres participants occupant un rôle d'observateur. Nous avons expliqué les critères d'observation : durée du parcours, plots touchés, sorties du tracé, hésitations, oublis de consigne, durée de regard détourné et réaction à un événement visuel. Nous avons précisé que les mesures n'avaient aucune valeur de diagnostic individuel et qu'elles servaient uniquement au débat pédagogique.

\paragraph{Premier passage : parcours de référence sans distracteur.}
Un volontaire a reçu la consigne de réaliser normalement le slalom, de s'arrêter à l'endroit indiqué, de toucher l'objet prévu et de franchir la ligne d'arrivée. Le temps et les éventuelles erreurs ont servi de référence pour comparer les conditions suivantes. Nous avons commencé le débriefing par une question ouverte sur la difficulté ressentie.

\paragraph{Deuxième passage : appel professionnel.}
Le même parcours a été repris avec la simulation d'un appel imprévu, préparé par un animateur : \og Bonjour, c'est urgent. Le client attend. Pouvez-vous me confirmer l'adresse du prochain chantier ? \fg{} Les observateurs devaient relever si le participant s'arrêtait, cherchait à décrocher, ralentissait, oubliait une étape ou touchait un plot. Nous avons interrogé le groupe sur ce qui aurait pu se passer si l'événement s'était produit au volant d'un utilitaire.

\paragraph{Troisième passage : lecture d'un SMS.}
Nous avons introduit un message bref de type professionnel : \og Nouveau chantier : 12 rue des Lilas. Le client attend depuis 20 minutes. \fg{} Le volontaire devait lire l'information et la restituer, sans accélérer et avec la possibilité de s'arrêter immédiatement. L'observation portait sur le regard détourné, la trajectoire, les hésitations et les consignes oubliées. Nous avons relié cet essai au cumul des dimensions visuelle, manuelle et cognitive.

\paragraph{Quatrième passage : double tâche et événement visuel.}
Le parcours a été associé à une information à écouter ou à lire ainsi qu'à un signal visuel simple : une balle en mousse présentée ou envoyée très doucement, sans viser la tête. Les volontaires étaient prévenus de la possibilité de cet événement, mais pas de son instant exact. Attraper la balle restait facultatif ; il était interdit de courir ou de se jeter pour l'atteindre. L'objectif était d'observer la disponibilité attentionnelle, sans mettre les stagiaires en situation de chute.

\paragraph{Cinquième passage : sollicitations cumulées, si les conditions le permettaient.}
Un dernier essai facultatif associait lecture d'une information, question orale et événement visuel. Nous avons utilisé ce cas comme analogie des situations professionnelles dans lesquelles le GPS, le client, le responsable et le retard sollicitent simultanément le conducteur. \emph{Le maintien de ce passage dans le récit définitif dépendra de sa réalisation effective le jour de l'action.}

\paragraph{Débriefing et transfert professionnel.}
Nous avons organisé la mise en commun des observations en distinguant les différents passages, sans présumer d'une augmentation systématique des temps ou des erreurs. Quatre questions ont structuré les échanges : qu'est-ce qui a changé ? Aviez-vous conscience des difficultés éventuelles ? Quelle sollicitation vous a le plus perturbés ? Quelle décision adopter avant de répondre à un appel ou à un message au volant ?

La comparaison était volontairement présentée comme une \textbf{analogie pédagogique} : toucher un plot dans un environnement sécurisé n'a pas la gravité d'une collision avec un piéton ou un cycliste. Un parcours à pied ne permet pas de quantifier le risque d'accident en voiture ; il aide à prendre conscience de la concurrence entre les tâches et à discuter des décisions protectrices.

\subsection{Organisation de la co-animation et sécurité}
Nous avons préparé une répartition fonctionnelle des responsabilités : un intervenant présentait les consignes et conduisait le questionnement, tandis que l'autre gérait le matériel, le temps et la sécurisation du passage. Les rôles pouvaient être inversés lors des séquences. La zone extérieure devait être dégagée, isolée de toute circulation et exempte d'obstacles dangereux. Les déplacements se faisaient uniquement en marchant. Les messages étaient fictifs, les observations anonymes et chacun pouvait refuser une sollicitation ou interrompre l'exercice.

La transition entre la salle et l'extérieur, ainsi que le retour collectif, avaient été pris en compte dans le minutage du tableau synoptique de deux heures. Quelques volontaires réalisaient les essais successifs pendant que les autres notaient des observations ; nous ne prévoyions pas cinq passages pour chacun des quinze stagiaires.

\section{Évaluation des acquis à chaud}
Nous avons terminé le scénario par un quiz oral inspiré de la diapositive « Testez vos connaissances ». Quatre questions portaient sur les types de distraction, la sanction du téléphone en main, la distance parcourue en cinq secondes à 90 km/h et l'importance du défaut d'attention dans les accidents. Nous avons également prévu une question d'application : \og Un client appelle alors que vous êtes au volant pour vous annoncer une urgence : que faites-vous ? \fg{}

L'évaluation devait vérifier non seulement la mémorisation d'une donnée, mais surtout la capacité à proposer une réponse sûre : ne pas traiter la sollicitation pendant la conduite, rechercher un lieu autorisé où stationner et rappeler à l'arrêt. Les réponses des observateurs au parcours étaient considérées au même titre que celles des volontaires. Aucune note individuelle n'était recherchée.

\textbf{Données à reporter après l'action :} nombre de participants présents et évalués, questions réellement posées, répartition des réponses, erreurs fréquentes et éventuelles reformulations. Aucune progression mesurée ni réussite collective n'est affirmée à ce stade.

\section{Questionnaire de satisfaction}
Nous avons préparé un questionnaire anonyme distinct du quiz, portant sur la clarté des explications, l'intérêt des deux vidéos, l'utilité du parcours pédestre, la pertinence des situations de chantier, le temps laissé aux échanges et les comportements que les participants envisagent de modifier. Une question ouverte permettait de proposer des améliorations.

\textbf{Données à reporter après l'action :} questionnaire effectivement administré, nombre de retours, synthèse des appréciations et propositions recueillies. Les intentions exprimées ne seront pas assimilées à un changement durable du comportement routier.

\section{Vérification du compte rendu après l'intervention}
\texteguide{Au retour du 4 décembre : vérifier le contenu réellement projeté (y compris la lecture effective des vidéos), les durées et l'effectif ; supprimer tout passage non réalisé ; renseigner les rôles effectifs du binôme, les observations du parcours, les résultats du quiz et la satisfaction ; remplacer le présent encadré provisoire par un compte rendu exclusivement factuel.}
